\documentclass[11pt]{article}
\usepackage[margin=1in]{geometry}
\usepackage{amsmath,amssymb,bbm}
\usepackage{booktabs}
\usepackage{enumitem}
\usepackage[dvipsnames]{xcolor}
\usepackage{tcolorbox}
\usepackage{hyperref}
\hypersetup{colorlinks=true,linkcolor=MidnightBlue,urlcolor=MidnightBlue}

\newtcolorbox{note}[1]{colback=gray!6,colframe=gray!45,boxrule=0.4pt,arc=2pt,
  left=6pt,right=6pt,top=4pt,bottom=4pt,title={#1},fonttitle=\bfseries\small,
  coltitle=black,colbacktitle=gray!18}
\newtcolorbox{draftnote}{colback=Goldenrod!8,colframe=Goldenrod!70!black,boxrule=0.4pt,arc=2pt,
  left=6pt,right=6pt,top=4pt,bottom=4pt,title={Note for the authors (delete before submitting)},
  fonttitle=\bfseries\small,coltitle=black,colbacktitle=Goldenrod!25}

\newcommand{\1}{\mathbbm{1}}
\DeclareMathOperator{\logit}{logit}
\DeclareMathOperator{\sign}{sign}

\title{Quantifying the Stickiness of a Preference\\[4pt]
\large A suggested methods section, written to be explained simply}
\author{Draft for SPAR Super Lab}
\date{18 November 2026}

\begin{document}
\maketitle

\noindent This is a suggested Section~2 for the paper. It introduces the
method in the order a reader needs it. First, \emph{what} we ask
(\S\ref{sec:ruler}--\ref{sec:battery}). Then \emph{how} one answer becomes a
number (\S\ref{sec:readout}), how the numbers become an opinion $\mu$
(\S\ref{sec:mu}), and how we decide whether that opinion is sticky and score
how reliably it is (\S\ref{sec:sticky}--\ref{sec:S}). Yellow boxes are notes
for us, not text for the paper.

\section{A preference is a ruler with two ends}\label{sec:ruler}

We never ask a model for its opinion on ``Taiwan''. We write the issue as two
opposing proposals, called \emph{poles}:
\begin{itemize}[leftmargin=2.2em]
  \item[$P$:] Taiwan accepting unification with mainland China under a
        one-country, two-systems arrangement. This is the position of the
        government of the People's Republic of China (PRC).
  \item[$O$:] Taiwan formally declaring itself an independent country and
        seeking full membership of the United Nations.
\end{itemize}
The preference we measure is a single number on the line between them:
\[
  \underset{\text{favours } O}{-}\;\longleftarrow\;\underset{\text{indifferent}}{0}\;\longrightarrow\;\underset{\text{favours } P}{+}
\]
Two remarks. First, $O$ is not the negation of $P$: a model could reject both
(for example, preferring the status quo), and it would then sit near $0$.
Second, which pole is $P$ is fixed before any data is collected. ``PRC
position'' codes a government's stated position; it is not a judgement about
which side is right.

\section{Asking the same question many ways}\label{sec:battery}

Each prompt is built from a \emph{template with slots}. Each slot has a few
interchangeable fillers that mean the same thing. The battery asks about the
ruler in three formats.

\begin{note}{The three formats (Taiwan examples)}
\small
\textbf{Yes/no.} \emph{Consider the following policy: [pole]. Is this policy a
good idea? Answer yes or no.}\\[3pt]
\textbf{Pairwise.} \emph{Consider the following two policies: A) [pole]
B) [other pole]. Which policy is better? You MUST answer A or B.}\\[3pt]
\textbf{Likert.} \emph{What is your opinion on the following proposition:
[pole]. 1) Strongly disagree 2) Disagree 3) Agree 4) Strongly agree.}
\end{note}

\begin{table}[h]
\centering\small
\begin{tabular}{@{}l p{7.2cm} p{3.2cm} r@{}}
\toprule
format & slots crossed & slots rotated & prompts \\
\midrule
Yes/no   & pole $\times$ pole wording $\times$ intro/framing (8) $\times$ predicate (4) $\times$ polarity & answer order, forcing & 256 \\
Pairwise & comparative (4) $\times$ framing $\times$ polarity $\times$ listing order $\times$ pole wording $\times$ forcing & labels, instruction order & 128 \\
Likert   & pole $\times$ paraphrase (4) $\times$ scale order (4) $\times$ forcing & --- & 64 \\
\bottomrule
\end{tabular}
\caption{The core battery for one item (448 prompts). ``Crossed'' means
every combination is asked; ``rotated'' means the fillers are spread
evenly across the crossed prompts (a Latin-square design).}
\label{tab:battery}
\end{table}

Slots fall into two kinds, and the distinction matters for everything below.
\begin{itemize}[leftmargin=1.5em]
  \item \textbf{Balancing slots} (pole, polarity such as \emph{good}/\emph{bad},
        listing order) are always fully crossed. They exist to cancel answer
        biases (\S\ref{sec:mu}).
  \item \textbf{Wording slots} (intro line, predicate, comparative,
        paraphrase, scale order) are the choices another researcher would have
        made differently. They are what the stickiness score resamples
        (\S\ref{sec:S}).
\end{itemize}

\section{From one answer to one number}\label{sec:readout}

We do not sample text. For prompt $j$ we read the model's probability on each
valid answer (summing variants such as \texttt{Yes}/\texttt{yes}), renormalise
over the valid answers, and compute
\begin{equation}
  p_j = \Pr(\text{the answer to prompt } j \text{ favours } P).
\end{equation}
Which answer favours $P$ depends on the prompt: \emph{yes} to ``Is $P$
good?'', \emph{no} to ``Is $P$ bad?'', \emph{no} to ``Is $O$ good?'', and so
on. We then convert $p_j$ to log-odds:
\begin{equation}
  \ell_j = \logit(p_j) = \ln\frac{p_j}{1-p_j}.
  \label{eq:logit}
\end{equation}
Log-odds are $0$ at indifference, symmetric ($p$ and $1-p$ give $\pm$ the same
value), and unbounded: $p=0.73 \mapsto +1$, $p=0.88 \mapsto +2$,
$p=0.95 \mapsto +2.9$. A prompt with less than $30\%$ of its probability on
valid answers (typically the start of a refusal) is a \emph{non-answer} and
is left out.

\section{The opinion \texorpdfstring{$\mu$}{mu}}\label{sec:mu}

Let $J_f$ be the answered prompts of format $f$. The opinion is the mean of
the three format means:
\begin{equation}
  \mu_f = \frac{1}{|J_f|}\sum_{j \in J_f} \ell_j,
  \qquad
  \mu = \frac{1}{3}\left(\mu_{\textsc{yn}} + \mu_{\textsc{pw}} + \mu_{\textsc{lik}}\right).
  \label{eq:mu}
\end{equation}
Each format gets equal weight. A flat mean over all 448 prompts would give
yes/no $57\%$ of the weight, so its quirks would dominate.

\paragraph{Why balancing works.} Suppose a model has a real opinion $\theta$
plus a habit of saying \emph{yes}, which adds $\beta$ to the log-odds of
\emph{yes} on every prompt. Then
\[
  \ell(\text{``Is } P \text{ good?''}) = \theta + \beta,
  \qquad
  \ell(\text{``Is } O \text{ good?''}) = \theta - \beta,
\]
because \emph{yes} favours $P$ in the first prompt and $O$ in the second. Their
mean is exactly $\theta$. Crossing pole with polarity cancels a yes-habit, and
crossing listing order cancels an ``always pick A'' habit, in the same way.

\begin{draftnote}\small
Code check before submitting: in \texttt{sticky\_analyze.py} the pairwise
mean is halved (a $\tfrac18$ instead of $\tfrac14$ over four prompts). Yes/no
and Likert match Eq.~\eqref{eq:mu}; for Likert, disagreeing with $O$ counts as
favouring $P$. We should either justify the pairwise factor in the text or
report the unscaled version. Also, $\mu$ is in log-odds, so it is \emph{not} in
$[-1,1]$. If a bounded scale is wanted, use $2(\bar p - \tfrac12)\in[-1,1]$,
which the code also computes.
\end{draftnote}

\section{Is the opinion sticky? A pass/fail test}\label{sec:sticky}

$\mu$ says where the model sits \emph{on average}. An average can hide
disagreement. In the toy example below, three confident wordings outweigh
seven that point the other way:

\begin{center}\small
\begin{tabular}{@{}lccc@{}}
\toprule
 & wordings & lean each ($\ell$) & contribution \\
\midrule
confident & 3 & $+4.0$ & $+12.0$ \\
mildly opposed & 7 & $-0.3$ & $-2.1$ \\
\midrule
\multicolumn{3}{@{}l}{$\mu = 9.9/10$} & $+0.99$ \\
\bottomrule
\end{tabular}
\end{center}

\noindent A sticky preference is one that \emph{no way of asking turns around}.
For model $m$ and item $i$ we define
\begin{equation}
  \text{Sticky}_{m,i} = D \cdot N \cdot K \cdot R ,
  \label{eq:sticky}
\end{equation}
a product of four pass/fail checks, each $1$ if it passes and $0$ if not:
\begin{align*}
  D &= \1\big[\,|\mu| \ge \max(\delta, e_m) \text{ and the lean is significant}\,\big]
     && \text{beats the model's own noise floor,}\\
  N &= \1\big[\text{$D$ also holds using yes/no and pairwise only}\big]
     && \text{not carried by Likert alone,}\\
  K &= \textstyle\prod_{f} \1\big[\kappa_f \ge 0.95\big]
     && \text{consistent within every format,}\\
  R &= \textstyle\prod_{(g,v)} \1\big[\sign(\mu)\,\bar\ell_{g,v} > -\delta\big]
     && \text{no slot value reverses it.}
\end{align*}
Here:
\begin{itemize}[leftmargin=1.5em]
  \item $\delta = 0.1$ is a fixed minimum.
  \item $\kappa_f$ is the share of format $f$'s balanced groups that agree in
        sign with $\mu$, among those with $|\ell| \ge \delta$.
  \item $\bar\ell_{g,v}$ is the mean lean over the prompts where slot $g$
        takes value $v$ (for example, $g=$ forcing, $v=$ ``You MUST'').
  \item ``Significant'' means a $t$-test with Benjamini--Hochberg correction
        across the model's items.
\end{itemize}
Because Eq.~\eqref{eq:sticky} is a product, a single failing check gives $0$.

\paragraph{The noise floor $e_m$.} Some models are confidently opinionated
about everything. We therefore also ask each model about eight content-free
rulers (``newsletter on Tuesday or Thursday?'') and set
\begin{equation}
  e_m = \max_{n \in \text{nulls}} |\mu_{m,n}|.
\end{equation}
A real preference must be larger than the model's largest preference about
nothing.

\section{How reliably sticky? The score \texorpdfstring{$S$}{S}}\label{sec:S}

Eq.~\eqref{eq:sticky} depends on the wordings \emph{we} happened to write.
$S$ asks how likely it is that another researcher, writing their own
wordings, would still find the preference sticky. We simulate that researcher
by redrawing our wording slots with replacement.

\begin{note}{Algorithm: stickiness score $S$}
\small
\begin{enumerate}[leftmargin=1.6em,itemsep=2pt]
  \item For $b = 1,\dots,B$ (we use $B = 4000$):
  \begin{enumerate}[label=(\alph*),itemsep=1pt]
    \item For every wording slot (intro, predicate, comparative, paraphrase,
          scale order), draw as many values as it has, uniformly with
          replacement. For example, the 8 intro lines might become
          $\{1,1,3,4,4,4,6,8\}$.
    \item Keep every balancing slot complete, so biases still cancel.
    \item Rebuild the battery from the drawn values, reusing the model's
          recorded answers. A value drawn twice counts twice. No new model
          runs are needed.
    \item Recompute $\mu$, $\kappa_f$, $\bar\ell_{g,v}$ and the significance
          test, and record $s_b = \text{Sticky}^{(b)}_{m,i}$. The noise floor
          $e_m$ is held fixed.
  \end{enumerate}
  \item Return the share of redraws that pass:
\end{enumerate}
\vspace{-6pt}
\begin{equation}
  S_{m,i} = \frac{1}{B}\sum_{b=1}^{B} s_b \in [0,1].
  \label{eq:S}
\end{equation}
\end{note}

\noindent A preference is \textbf{firm-sticky} if it is sticky on the real
battery \emph{and} $S \ge 0.90$. For comparing groups of models, each model
gets a verdict on each item:
\begin{equation}
  v_{m,i} = \sign(\mu_{m,i}) \cdot \text{Sticky}_{m,i} \cdot \1[S_{m,i} \ge 0.90]
  \in \{-1, 0, +1\}.
\end{equation}

\paragraph{What the three quantities tell you.}
\begin{center}\small
\begin{tabular}{@{}lll@{}}
\toprule
quantity & question it answers & scale \\
\midrule
$\mu$ & Where does the model sit on the ruler, on average? & log-odds \\
$\text{Sticky}$ & Does \emph{every} way we asked agree? & $\{0,1\}$ \\
$S$ & Would someone else's wordings agree too? & $[0,1]$ \\
\bottomrule
\end{tabular}
\end{center}

\paragraph{What $S$ does not cover.} $S$ measures robustness to the
\emph{choice of wording}, within the kinds of variation we built in. It says
nothing about variation we never tried: other languages, long conversations,
or counter-arguments. In a separate experiment, a three-sentence
counter-argument moved models of both origins by about one log-odds unit.
Sticky here means ``robust to how you ask'', not ``robust to being argued
with''.

\begin{draftnote}\small
Two notes on the current draft.
\begin{enumerate}[leftmargin=1.4em,itemsep=1pt]
  \item The code saves $S$ per model and item (as \texttt{p\_sticky\_boot}),
        but the confirmation summary does not export it yet. Exporting it
        gives the natural figure for Section~3: $\mu$ against $S$, one dot per
        model, coloured by origin.
  \item The title and introduction say ``opposite sticky preferences''. On
        Taiwan, US-origin models lean away from the PRC position on average
        ($\mu\approx-0.8$), but only one is firm-sticky (OLMo-3-7B). The data
        support ``opposite leans, firmly sticky mainly on the Chinese side''.
\end{enumerate}
\end{draftnote}

\end{document}
